\cvsection{Expériences professionnelles}

\begin{cventries}
  \cventry
    {Stagiaire ingénieur de recherche en IA}
    {INRIA}
    {Bordeaux, France}
    {Mars 2026 -- Septembre 2026}
    {
      \begin{cvitems}
        \item {Analyse d’un besoin de recherche documentaire et développement en autonomie d’une application fonctionnelle intégrant un agent IA pour aider les chercheurs à trouver des articles scientifiques.}
        \item {Développement de l’agent multi-étapes sans framework d’orchestration, avec intégration de l’API OpenAI, conception des prompts et gestion des appels d’outils, de la mémoire et du contexte.}
        \item {Développement d’API REST en Python/FastAPI documentées avec OpenAPI, stockage dans MongoDB et création d’une interface React/Next.js pour rendre l’agent et ses outils accessibles aux chercheurs.}
        \item {Mise en œuvre d’un pipeline RAG avec chunking, embeddings, indexation dans ChromaDB et recherche sémantique. Automatisation de la collecte de métadonnées d’articles en ligne.}
        \item {Déploiement de LLM avec vLLM sur cluster HPC équipé de GPU NVIDIA et gestion de modèles locaux avec LM Studio sur Mac Studio.}
        \item {Définition d’un protocole d’évaluation et comparaison de LLM et de quantifications selon la réussite de la recherche, la latence et la taille du contexte.}
        \item {Conteneurisation de composants avec Docker et mise en place d’un pipeline CI/CD sur GitHub.}
      \end{cvitems}
    }

  \cventry
    {Stagiaire ingénieur de recherche en apprentissage par renforcement}
    {Université Karunya}
    {Coimbatore, Inde}
    {Juin 2025 -- Septembre 2025}
    {
      \begin{cvitems}
        \item {Développement d’un modèle d’apprentissage par renforcement pour la navigation autonome d’un robot hospitalier, puis déploiement et tests sur Raspberry Pi.}
      \end{cvitems}
    }
\end{cventries}
